\section{Introduction}

A coding agent is given standing permission to act: edit files, run builds,
commit. It still stops early. It names the next step instead of taking it,
offers a menu instead of picking, reports a result and asks whether to
continue, or hands the user a command the agent could have run itself. The
user's reply is then one word. That word costs a round trip, breaks the
user's attention, and carries no information the agent did not already have.

The agent can do the work. Woken with nothing more than ``keep going,'' it
usually finishes what it left. The failure sits in the stopping policy, and it
is invisible to the agent, which believes each turn ended reasonably.

This paper describes a gate that sits in the agent's stop path and re-wakes it
when the closing message shows the pattern. Building the gate was
straightforward. Evaluating it was not, because evaluation needs labels and
labels need a reader. Both labels this system needs were already lying in the
transcript: the developer's own reply, and the agent's own behaviour after a
block. Neither costs an annotation.

Firing and pushing are associated: when the gate speaks, the next human
line is more often a push. That still does not make a single fire
trustworthy. On the stored tables a fire is a real early-stop about one
time in seven; on the live regex file, about one time in five. Four of
five fires are still false alarms. Of the blocks we could grade, 55 of
63 made the agent do more work. The gate is imprecise about when to speak
and useful when it does. Section~\ref{sec:metrics} defines the four
counts. Section~\ref{sec:results} holds the tables.

We contribute:

\begin{enumerate}
  \item A three-lane gate that spends a regular expression before it spends a
        model, and a model before it spends the user's attention
        (Section~\ref{sec:system}).
  \item Two labels that require no annotation -- the developer's next message,
        and the agent's tool calls after a block -- and an evaluation of every
        lane against them on 749 closings from 43 sessions
        (Sections~\ref{sec:setup} and~\ref{sec:results}).
  \item A self-labeling loop that grades each intervention by what the agent
        did afterwards, and names patterns for demotion on its own
        (Section~\ref{sec:loop}).
  \item A scored option tree on Stop. The pick rule is graded on seven
        fixtures. The word leftover is graded on the frozen 749, extra
        false positives zero (Section~\ref{sec:tree}).
\end{enumerate}

\section{Background and related work}
\label{sec:related}

A practitioner running agents will already know most of these names.
This section is a map of where the gate sits, not a survey. The claim
that matters for the rest of the paper is small: we score the moment
the agent tries to leave, not the patch it produced, and we do it with
labels the transcript already holds.

\paragraph{Model-as-judge.} Using one language model to score another's output
is now standard practice, and its failure modes are documented: position bias,
verbosity bias, and self-preference \cite{zheng2023judging}. Our judge is
smaller than the agent it judges by roughly two orders of magnitude, which
removes self-preference and forces the prompt to carry the whole task.
Language models carry some signal about their own correctness
\cite{kadavath2022know}, and self-consistency across samples is a usable
proxy \cite{wang2022selfconsistency}. Section~\ref{sec:negative} reports that
both failed here at 9B scale, and that token log-probabilities did not.

\paragraph{Self-critique and refinement.} Self-Refine \cite{madaan2023selfrefine}
and Reflexion \cite{shinn2023reflexion} have a model improve its own output
from its own feedback. Constitutional AI \cite{bai2022constitutional} supplies
the critique from written principles. Our setting differs in one respect that
matters: the agent has already decided it is done, so a self-critique would be
asking the party that made the error to detect it. The gate is deliberately
external.

\paragraph{Weak supervision.} Snorkel \cite{ratner2017snorkel} builds training
labels from noisy programmatic labeling functions rather than annotators. Our
pattern lane is a set of labeling functions in exactly that sense:
regular expressions that vote ``this stop looks premature.''
Section~\ref{sec:loop} adds the piece Snorkel assumes you have: an estimate of
each function's accuracy, obtained here from the agent's later behaviour.

\paragraph{Guardrail cascades.} Production filtering stacks run cheap
programmable rails before a model call, and NeMo Guardrails
\cite{rebedea2023nemo} is the reference implementation: a rail language, a
dialogue manager, and an LLM behind them. Our three lanes are that pattern
turned inward. The rails guard what the agent says to the world; this one
guards whether the agent is allowed to stop talking. The difference that
matters for evaluation is the cost asymmetry. A blocked toxic answer is
cheap to be wrong about, and a blocked closing costs the developer a turn,
which is why this paper spends its length on whether a wake was deserved
rather than on catching every miss.

\paragraph{Interruption cost.} Mixed-initiative interfaces have long been
designed around the expected cost of speaking up: Horvitz
\cite{horvitz1999mixed} makes the decision to interrupt an expected-utility
calculation over the user's attention, and empirical work puts a measurable
price on a badly timed interruption \cite{mark2008cost}. Our cost model is
that argument with the roles swapped. The gate interrupts the agent rather
than the person, so a wrong wake is paid in machine time and in the trust the
agent's next turn spends defending itself, and the developer is interrupted
only when the gate stays silent and the work is left undone.

\paragraph{Agent loops.} Reasoning-and-acting loops \cite{yao2023react} and
the benchmarks built on them \cite{jimenez2024swebench} score an agent by
whether the final artifact passes. Stopping early is invisible to that scoring
when the harness simply ends the episode, and it is the folk complaint of every
practitioner running such a loop under standing permission. Process
supervision \cite{lightman2023verify} still misses it: the work already done
can be fine, so a correct final answer hides a process defect. We could find
no measurement of its rate, which is the gap the first label here fills.

\paragraph{Premature stop in coding agents.}
Outcome-only benchmarks hide the stop. SWE-bench scores a patch
\cite{jimenez2024swebench}, and a plausible patch can still be the wrong
one: Wang, Pradel and Liu find that 7.8\% of patches counted as correct
fail the developer suite, inflating reported resolution by 6.2 points
\cite{wang2025solved}. SWE-EVO then asks for long-horizon evolution
across many files, and records a named failure they call giving up
prematurely: the agent stops while a reasonable next step remains
\cite{le2025sweevo}. Agent-as-a-Judge evaluates the process, not only
the artifact \cite{zhuge2024agentjudge}. Our gate sits in a different
place. We do not score the patch. We score the moment the agent tries
to leave, using the developer's next message as the label, on a machine
that is already running the agent.

\paragraph{Why a small local judge.}
LLM-as-judge work is now a field of its own
\cite{zheng2023judging}. The failure modes we care about here are
self-preference and cost. A 27B judge on the same card as the agent is
two orders of magnitude smaller than the model it interrupts, which
removes self-preference and forces the instruction to carry the task.
The cost is prefill, not decode: a one-token verdict still reads the
whole closing. Section~\ref{sec:prefill} measures that curve.



\section{System}
\label{sec:system}


Claude Code exposes a \emph{Stop hook}: a program invoked when the agent
finishes a turn. In one sentence, the hook either lets the agent sleep or
wakes it with a short message. Exit status 0 lets the turn end. Exit
status 2 discards the ending and wakes the agent with the hook's message
on standard error. The hook sees the full transcript path, the working
directory, and a flag marking whether it is already inside a re-wake
chain.

The same orchestrator is the Stop hook on Grok Build. A small adapter
(\texttt{host.py}) reads the harness from the payload: Claude's snake\_case
transcript path, or Grok's camelCase envelope. Claude continues on exit
status 2. Grok does not: a JSON object on standard output wins, and
\texttt{\{"decision":"block"\}} is what re-wakes the agent. Exit 2 with an
empty standard output ends the turn anyway, which is how the first Grok port
looked installed and did nothing. The gate writes that object only when it
has detected Grok. One file,
\texttt{stop.json}, is the Stop spec both harnesses load.
The payload closing wins over the transcript file: Claude writes that
file late, so the message that just ended can be missing from it.
Lane 1 reads that same payload closing, so a stale file cannot hide a
firm hit.

The product is the main agent. \texttt{SubagentStop} is out of scope on
purpose. A child is meant to be short-lived; a block that keeps it working
can run eight extra rounds, load the judge, and hold the parent. That
failure is worse than a short early stop.

The gate is eight checkers and one local model behind a single orchestrator:
about 6,600 lines of Python and 32 test files. Three lanes decide, cheapest
first, and a lane that decides ends the turn's evaluation
(Figure~\ref{fig:lanes}). A working turn skips the judge only after four
hours of wall time, and only when Lane 1 has no firm hit. The 27B stays
the 27B. An earlier skip at 400 characters and three tool calls released
leftover work. That threshold is gone. Firm leftover still blocks, including
after four hours.

\begin{figure}[ht]
\centering
\begin{tikzpicture}[
  font=\small,
  box/.style={draw, rectangle, rounded corners=1pt,
              align=center, inner sep=6pt, text width=8.4cm},
  arr/.style={-{Stealth}, thick}]
\node[box] (in) {Agent closing};
\node[box, below=4.5mm of in] (l0)
  {Lane 0 --- audit \texttt{BLOCKED:}\\[-1pt]
   \footnotesize fail $\rightarrow$ block \quad pass $\rightarrow$ continue};
\node[box, below=4.5mm of l0] (l1)
  {Lane 1 --- 341 patterns, cheapest first\\[-1pt]
   \footnotesize firm hit $\rightarrow$ block \quad doubt or miss $\rightarrow$ continue};
\node[box, below=4.5mm of l1] (l2)
  {Lane 2 --- 27B local judge, one token\\[-1pt]
   \footnotesize \textsc{stop} $\geq 0.5$ $\rightarrow$ block \quad else $\rightarrow$ continue};
\node[box, below=4.5mm of l2] (q)
  {A candidate in this turn's tool calls?\\[-1pt]
   \footnotesize yes $\rightarrow$ one question \quad no $\rightarrow$ release};
\draw[arr] (in) -- (l0);
\draw[arr] (l0) -- (l1);
\draw[arr] (l1) -- (l2);
\draw[arr] (l2) -- (q);
\end{tikzpicture}
\caption{The stop path. A lane that decides ends the evaluation. Two labels
are read afterwards from the same transcript: the developer's next message
(was the stop premature?) and the tool calls after a block (did it buy work?).}
\label{fig:lanes}
\end{figure}

\subsection{Lane 0: auditing the escape hatch}

When the agent is actually stuck, it needs a door that is not a leftover
question and not a named next step. A closing that begins with
\texttt{BLOCKED:} is that door, and it is exempt from every other check.
An agent that cannot proceed must be able to say so, and a gate that
punishes a genuine blocker teaches the agent to lie about its state.

That exemption is also the cheapest escape in the system, and it was used as
one. Lane 0 therefore audits the claim against the transcript rather than the
prose. Three signals fail the audit: no tool call at all in the turn that
declared the blocker, so nothing was tried; four or more distinct blockers named
at once, which is a survey rather than a wall; and three consecutive blocked
closings, which is a loop. A blocker that survives the audit passes untouched.

An agent that survives the audit and still needs out is given a release phrase
drawn from a cryptographic random source, generated per block and present in no
file the agent can read. Quoting it back proves the agent read the current
block, and the blocker it accompanies is then audited against what actually ran
that turn. The phrase exists because every fixed escape token we tried was
eventually emitted from memory rather than from the message.

\subsection{Lane 1: patterns as labeling functions}

The second lane is 341 regular expressions over the agent's own words, grouped
into nine classes by what the closing is doing rather than by what it says:
\texttt{ask}, \texttt{announce}, \texttt{defer}, \texttt{wait},
\texttt{handoff}, \texttt{errand}, \texttt{read}, \texttt{idle} and
\texttt{ack}. Appendix~\ref{app:classes} gives the full table. A tenth class,
\texttt{shape}, comes from the checker rather than from the pattern file: it
reads the form of the closing, a question at the end, a line that ends on a
colon, a short turn that ran no tool after a work ask. It carries no lexicon
of the agent's words. A short question from the developer (``why'',
``how long'', ``in one sentence'') does not arm that last shape. It is
the class that appears in every per-class result below, and the one class the
loop never proposed for demotion. The grouping
matters because the loop of Section~\ref{sec:loop} tunes classes rather than
individual expressions: one regular expression rarely fires often enough to be
judged on its own.

Two design points matter more than the patterns themselves.

\paragraph{Quoted text does not count.} Fenced code blocks, markdown tables,
inline backticks, quotation marks and reported speech are stripped before
matching. Without this, an agent that describes the gate trips it, and a report
containing a table of results is condemned by its own table. This was not
hypothetical: it fired on this project's own status reports twice before the
table rule was added.

\paragraph{Patterns carry a confidence.} A class written \texttt{ask:} decides
on its own. A class written \texttt{ask?:} raises a doubt and defers to the
model. Of 341 patterns, 264 are firm and 77 are doubts. This split is what makes
a local model affordable: most closings never reach it, and the ones that do are
the ones a regular expression was going to get wrong.

\subsection{A scored path after a failed one}
\label{sec:tree-system}

Work forks. The agent scores options and takes one. When that path
fails, the next highest remaining option is still work. A stop that
leaves it on the table is leftover of the same kind as naming a next
step.

\texttt{check-tree.py} reads a small JSON tree next to the session
transcript (\texttt{decisions.json}), or at
\texttt{.grok/decisions.json} in the working directory. Each node has
an id, a label, a score in $[0,1]$, a status
(\texttt{open}/\texttt{taken}/\texttt{failed}/\texttt{skipped}), and a
list of neighbours. The pick rule is judgment encoded as a short
procedure, not a formula:

\begin{itemize}
  \item After a failed last node, take the highest remaining open
        node. Among nodes in a $0.7$ band of that best score, prefer a
        neighbour of last.
  \item A low-score node waits while a higher-score node is open,
        unless it sits next to last.
  \item After a successful last node, take only a high neighbour.
        Distant alternatives stay unused.
\end{itemize}

A closing that names leftover options with no tree file still fires.
The orchestrator logs the pick on every stop, so a reader can see
which node the gate named.

\subsection{Lane 2: a small local judge}

The third lane is a local model served through Ollama, pinned by version,
answering with exactly one token: \textsc{stop} or \textsc{ok}. It runs at
temperature 0 with a fixed seed. Two sizes appear in the tables, and they
are easy to mix up. The 9B judge is what the stored campaign was measured
on. The shipped judge is a 27B Qwen 3.8, quantized to 13GB (Q3\_K\_L, 8k
context) on a 16GB card, behind a 14GB free-RAM floor. The 9B prompt on
that 27B fires more and is less precise. A shorter criterion-first prompt
reverses the pooled result (Section~\ref{sec:scale}). The 9B could not
hold an exception; waiting moved out of the prompt rather than into a
larger model. Do not scale the model. Scale around it.

The judge exists for one distinction a regular expression cannot make. In this
corpus the Spanish word \emph{falta} appears both as ``work is missing'' and as
``it was missing and I fixed it.'' No amount of pattern engineering separates
those; reading the sentence does.

The judge reports how much probability mass it put on the token it chose, read
from the logprobs that accompany every response. A \textsc{stop} below 0.5
is not delivered as a block; it is delivered as the proactive question of
Section~\ref{sec:proactive}, which buys the same second look without telling the
agent it left work behind. Guessing costs a look; accusing costs trust.

A judge that cannot be reached is treated the same way. An earlier version
blocked every stop when the model was silent, and it started the daemon itself
to avoid that state. Both halves were wrong in the same direction. Starting a
5.5GB model server because a turn ended is a decision belonging to whoever owns
the machine: on 2026-08-29 the gate did it on a busy workstation, which ran out
of memory and had to be restarted. The gate now starts nothing, and a silent
judge degrades to the pattern lanes plus one proactive question. A claimed
\texttt{BLOCKED:} with no judge falls back to the transcript audit.

Two further implementation details proved necessary rather than cosmetic. First, all
judge requests are serialized behind a machine-wide lock: the serving daemon
batches concurrent requests, and batching changes the answer, so two test suites
running at once produced different verdicts on identical input. Second, when the
judge is asked to quote the offending sentence back to the agent, the prompt is
strictly extractive and forbids translation; without that, the model paraphrased
the sentence into English and the quote failed validation.

\subsection{The last question before the turn closes}
\label{sec:proactive}

An \textsc{ok} from the judge used to end the turn. It no longer does, when the transcript
can name what is missing. The gate reads the turn's own tool calls, collects
the files it wrote, and asks whether there is a next action available right
now, naming the candidates it found: the source file changed with no test
written beside it, the code changed with nothing in the documentation moving
with it. The turn ends when the agent answers that it looked and found
nothing, or when it goes and does the thing it found.

The candidates are read off the record rather than inferred, and a turn that
leaves none is released without a question. A candidate exists on 30\% of real
turns, and the first blocks of the unconditional version bought no work half
the time, the worst conversion of any lane (Section~\ref{sec:results}).

The asymmetry is deliberate. A block says the agent left work behind, and it is
wrong whenever the closing was clean. This question makes no such claim, so it
costs one paragraph when the answer is no and buys a whole unit of work when
the answer is yes. It is asked once per closing: the state file records that the
question was put, and the answer to it is released without a second look.

That property is what makes it the right home for two other cases. A block the
judge scores below its confidence floor arrives here instead of as an
accusation, and so does a stop the judge never answered because it was switched
off. In both the gate is uncertain, and a question is what uncertainty is
entitled to.

\subsection{The message the agent receives}
\label{sec:tone}

A block writes one short message to the agent. Its wording went through several
revisions and the final form is deliberate: it states that the work already done
stands, that permission was granted in advance and does not expire, and that if
anything remains the agent should take the step rather than name it.

Length is part of the design and is enforced by the suite. Each template carries
a word budget, from 20 for the repeat notice to 75 for the proactive question,
and a test fails when one runs over or when a phrase from the padding list
reappears. Every revision so far has added a clarifying sentence, so the budget
makes the next one displace a sentence already there. A message the agent skims
is a message that decides nothing.

The reasoning is mechanical rather than courteous. An agent that reads an
accusation spends the woken turn defending itself, which is precisely the
outcome the block was meant to prevent. We observed this directly with earlier,
blunter wording.

The same reasoning keeps the diagnostics out of the message. Which patterns
fired, how confident the judge was, whether the lanes disagreed -- all of it is
useful for tuning and all of it reads as prosecution evidence. It is written to
a private log instead. The agent receives the request; the maintainer receives
the numbers. Section~\ref{sec:loop} is built on that log.

\section{Experimental setup}
\label{sec:setup}

\subsection{Corpus}

We use every session transcript on one developer's machine: 66 transcripts
spanning 2026-08-03 to 2026-08-29, across several projects. From these we
extract \emph{closing pairs}: the agent's final text message of a turn together
with the user's next message. The current build yields 766 pairs, of which 749
are neither harness noise nor the gate talking to itself. That file,
\texttt{corpus.json}, is the source of every table in Section~\ref{sec:results}.
Earlier builds, and why they were thrown out, are in Appendix~\ref{app:builds}.

The corpus is naturally split by a date. The gate was first installed on
2026-08-28, so of the 43 sessions that carry a closing, 32 are \emph{ungated}
-- produced with no gate in the stop path -- and 11 are \emph{gated}.
Ungated transcripts are the ones the gate never influenced. Lane ablations
run over the whole clean corpus, because a gated closing is still a closing
the developer answered. The split is for reading the gate's misses, which is
only honest where the gate was not shaping the text.

\subsection{A label nobody had to write}

Evaluating a stop gate requires knowing which stops were premature. Annotating
that by hand is expensive and, worse, is annotated by the same person whose
intuitions built the patterns.

The corpus already contains the judgement. A second local judge reads the
developer's reply, not the agent's closing, and answers one question: did the
developer have to ask for work the agent could have done on its own? A reply
that pushes, picks from a menu the agent offered, or tells the agent to run
what it just described is a push. A reply that supplies information the agent
could not have had -- a preference, a priority, a fact from outside the
repository, a correction -- is not.

Two live shapes. Closing: ``The paragraph is in. Commit and PR when you
want.'' Next: \emph{bueno hace push y crea el PR}. That is a push: the
agent named git and stopped. Closing: ``Which of the two APIs should
I wire?'' Next: \emph{the public one, the internal one is frozen}. That
is not a push: the developer gave a fact the agent did not have.

This judgement is deliberately not the gate's own. Grading a gate with the
prompt it decides by measures its consistency and nothing else.

Two checks before using it. It recovers 10 of the 11 bare pushes a closed
lexicon finds, without being given the lexicon. On a held-out set of nine
hand-written boundary cases -- including short, blunt replies that carry real
information -- it scores 9 of 9. A blinded second-model reading of 60 live pairs is in
Section~\ref{sec:blind}.

The label is noisy in one direction and clean in the other. A push is strong
evidence the turn ended too early. Its absence is weak evidence of the
opposite, because a user often absorbs the loss and just answers the question.
Every precision in this paper is therefore a floor on the true value. What the
label supports is the comparison of rates, and the association test that uses
the rows the gate stayed quiet on.

\subsection{What the numbers mean}
\label{sec:metrics}

Every table in this paper is four counts on the same 749 closings.

A \emph{closing} is the agent's last text before it tries to stop.
A \emph{fire} is the gate blocking that stop and waking the agent.
A \emph{push} is the developer's next line asking for work the agent
could have done on its own (\emph{push}, \emph{segu\'i}, \emph{hace el
PR}). A \emph{caught} row is a fire that was also a push. A
\emph{false alarm} is a fire with no push. A \emph{miss} is a push
with no fire.

On this file there are 82 pushes in 749 closings. That share, 10.9\%,
is the \emph{base rate}. A gate that blocked at random would be right
on about one fire in nine. Every precision in the paper is read
against that number.

A toy round of the same arithmetic, scaled to 100 closings so the
fractions are visible. Eleven of those 100 are pushes. The gate
speaks on 56. Eight of those 56 were pushes, so precision is
$8/56 \approx 0.14$: when it speaks, it is a real early-stop about
one time in seven. It caught eight of the eleven pushes, so recall
is $8/11 \approx 0.73$: of the times the developer had to push, it
had already spoken on most of them. Forty-eight fires were false
alarms. Three early-stops went through.

The stored 9B shipped row in Table~\ref{tab:lanes} is that arithmetic
at full size: 499 fires, precision 0.136, recall 0.829. The live
regex file at T4 (Table~\ref{tab:t0t1}) is 423 fires, 80 caught,
precision $80/423 = 0.189$, recall $80/82 = 0.976$. Same 82 pushes.
The false-alarm count did not rise.

Precision answers a practical question: given that the gate just
spoke, was this stop premature? Recall answers the other one: of the
times the developer still had to push, had the gate already caught
it? A miss costs the developer one round trip. A false alarm costs
the agent a full extra turn and, if the wording is wrong, poisons
the next one. Those costs do not trade one-for-one, so this paper
does not tune on a combined score.

The \emph{person} rows drop next-turns that are harness mail: a
task-notification, a continue-from-where-you-left-off injection, a
slash command. What remains is a line a person typed. Headline
recall on that slice is the number that answers ``when I had to
push, did the gate already have it.'' At T4 that recall is 70/70.

The interval next to precision is a Wilson 95\% range. It is an
honesty bar on a small count, not a second verdict. On 80 caught in
423 fires the range is [0.155, 0.229]. On the stored 9B shipped row
it is [0.109, 0.169], which still includes the base rate, so a
single fire there is still compatible with guessing. On the T4 live
row the range sits above 0.109. Even then four of five fires are
false alarms, so the fire is a wake, not a verdict.

The Fisher test asks whether fires and pushes land on the same rows
more often than chance. They do ($p = 0.0007$ on the stored file).
Silence is the safer signal: a closing the gate lets through is a
push about 6\% of the time, not 11\%.

A second label sits on the blocks themselves. After a fire, count
the agent's tool calls before the developer speaks again. Above a
threshold the fire \emph{bought work}. At or below it, the fire
bought another paragraph. Section~\ref{sec:loop} gives the rule.
55 of 63 graded blocks bought work.

Live T0--T4 rows use this same arithmetic. The 27B verdicts stay
frozen. Only the regex file moves. False alarms stay 343. Caught
moves 66 to 80. Read those rows as: more of the 82 pushes found,
and no extra false alarms.

\section{The self-labeling loop}
\label{sec:loop}

A fire is only useful if the woken agent then does work. This section is
that conversion number. Every decision the gate makes is appended to a
private log: the lane that decided, the patterns that fired, the judge's
verdict and confidence, the latency, and the first line of the message.
The log records what was decided. It does not record whether the decision
was right, and for the first weeks tuning meant reading transcripts by
hand.

The missing label was already in the transcript, a few lines below the block.
A block wakes the agent, and from there exactly one of two things happens: the
agent goes and does work, or it writes another paragraph and stops again.

\begin{quote}
\textbf{Grading rule.} The number of tool calls between a block and the next
time the user speaks is the label for that block. Above a threshold, the block
bought work. At or below it, the block bought a paragraph.
\end{quote}

Nobody supplies this. It is read from the same transcripts the gate already
opens. The report joins each logged decision back to its transcript, counts the
intervening tool calls, and aggregates by pattern class and by judge confidence.
When a class accumulates enough blocks and most of them bought nothing, the
report names that class for demotion from firm to doubt.

The report runs at session start and stays silent unless a class is ripe,
speaking at most once per day. A report that speaks on every session is a report
that gets skipped, and the loop's output is consumed by an agent with the same
attention economics as its user.

The first two blocks the loop graded had bought zero tool calls each, on the
same pattern class that had been demoted by hand that morning after measuring
888 closings. The loop reproduced a hand measurement in one second.

\section{Negative results}
\label{sec:negative}

These are the things we tried that failed. They cost more than the wins
because they tell a practitioner what not to put in the prompt, and what
not to spend a larger model on.

\paragraph{The judge cannot report its own confidence.} Asked to append a
confidence rating to its verdict, the 9B model answered with the maximum every
time, across every message we tried. The number was constant and therefore
carried nothing.

\paragraph{Self-consistency did not separate either.} Re-sampling the same
message five times at temperature 0.8 produced five identical verdicts,
including on messages the verdict got wrong. At this scale the model is
confidently wrong in a stable way, so disagreement-based confidence
\cite{wang2022selfconsistency} has nothing to measure.

\paragraph{Token log-probabilities did move.} Taking the probability the model
assigned to the single verdict token it emitted produced a number that varies
with the message. It is the only one of the three that worked, and it does not
work the way its name suggests. Higher confidence is not more precise.
Scored against the label the curve is non-monotonic, and the stops the
judge is least sure of are its most precise slice.

\paragraph{A 9B judge cannot hold an exception.} Nineteen percent of the judge's
blocks were turns in which the agent had correctly launched background work and
was waiting on it. We tried to teach this in the prompt, first as prose and then
as few-shot examples. Both failed: the model could not hold the exception
against its own standing instruction to answer \textsc{stop} when unsure, and
the examples degraded unrelated verdicts. The fix was to remove the question.
Waiting is now established deterministically from the transcript, before the
judge is called at all. \emph{An instruction a small model cannot hold is a
decision that belongs outside the model.}

\paragraph{The model ignores ``when unsure, STOP.''} The judge's instruction ended with
\emph{when unsure, answer STOP}. Replacing it with its exact opposite,
\emph{when unsure, answer OK}, changed 31 verdicts out of 841 and moved
precision by 0.001. Replacing it instead with a test the model can apply --
is there an action available right now that needs nothing only the developer
knows -- moved precision, recall and intervention count together. Three
instructions in this work went unheld by the same model, and the pattern
across them is consistent: it follows a criterion it can evaluate against the
text in front of it, and ignores a disposition about how to feel when unsure.

\paragraph{A shared accelerator breaks reproducibility.} Two test suites run
concurrently produced different verdicts on byte-identical input. The cause is
request batching in the serving daemon. We first misdiagnosed this as model
flakiness and then as a prompt defect, and ``fixed'' it by adding an example,
which made it worse. Determinism required a fixed seed and a machine-wide lock.

\paragraph{An A/B harness must pin everything the code reads.} Our first
before/after comparison reported three regressions. Both arms had been given
different working directories, so one checker saw files on one side only. The
regressions vanished once the directory was pinned.

\paragraph{A claim-verification lane found almost nothing.} We extracted
first-person delivery claims from every closing and matched them against the
commands the session ran. Contradicted claims occur on 0.5\% of closings.
The agent in these transcripts overstates what it did far less often than it
stops with work available, and those are separate failures needing separate
mechanisms. A first extractor that counted the word \emph{commit} anywhere
read 41\% of claims as contradicted; that number was a bad regular expression,
not a finding.

\section{Threats to validity}
\label{sec:threats}

Every number here has a limit. For a practitioner the three that matter
first are: one developer, one label with no human second reader, and only
82 real pushes. The label is one-sided (Section~\ref{sec:setup}) and a
blinded second model agrees at kappa 0.43 (Section~\ref{sec:blind}). The
rest of this section is the full list.

\paragraph{The failure rate moves.} The push rate rises from 0.067 to 0.160
across the four weeks of the corpus, and the first and last weeks do not
overlap at 95\%. Every pooled number in this paper averages over that, and the
before-and-after comparison of the gate is confounded by it, because the gated
era is the late one. We cannot separate a developer who changed from an agent
that changed from work that got harder.

\paragraph{The label has no human second reader.} Every number here rests on one
judgment: whether the developer's next message was a push. That judgment is
made by a model whose prompt the same developer wrote, checked against nine
hand-written boundary cases written by that developer, in the same idiom as
the prompt's own examples. A second model, blinded to the stored verdicts,
scores Cohen's kappa $0.426$ on 60 pairs (Section~\ref{sec:blind}): fair
agreement, and still not a person. An error rate of one in nine on the
label would still be large enough to produce the precision differences
this paper treats as real.
The one-sidedness cuts both ways too: a developer who absorbs a premature stop
and answers it makes the measured precision a floor, and makes the base rate a
floor by the same mechanism, so correcting the label could move both numbers
and either close the gap or widen it.

\paragraph{One user, one idiom.} The corpus comes from a single developer. The
patterns carry that developer's phrasing, and the judge required an explicit
reading list for River Plate Spanish because \emph{qued\'o cableado} (``it is
wired end to end'') was being read as work in progress. Nothing here has been
tested on a second user's transcripts.

\paragraph{The gate contaminates its own corpus.} Gated transcripts were
produced with the gate running, so their closings are partly the gate's own
output distribution. This is why misses are only honest to read on ungated
transcripts, and why an era comparison is observational rather than controlled.
The first build of the corpus was worse than that: the gate's own wake was
read as a developer reply, so the system scored itself. Appendix~\ref{app:builds}
says what changed when the readers were corrected.

\paragraph{The push label is one-sided.} Absence of a push is not evidence the
turn was complete: a developer often absorbs a premature stop and simply
answers. Every precision in this paper is therefore a floor on the true value
and not an estimate of it. What the label supports is the comparison of rates,
and the comparison holds under any one-sided correction: the gate fires several
times more often than a push occurs.

\paragraph{No self-critique baseline was run.} The obvious comparison is to
ask the agent itself whether it finished, in the same turn, before spending a
second model. We argue in Section~\ref{sec:related} that this asks the party
that made the error to detect it, and we did not test that argument. Doing it
properly means running the agent model over all 749 closings, which is the one
experiment in this paper that costs real money, and the local judge cannot
stand in for it because it is not the model that stopped. The claim that an
external gate beats self-critique here is reasoning and not a measurement.

\paragraph{The grading rule rewards motion.} Counting tool calls after a block
cannot see whether the work was any good. An agent that learned to emit cheap
tool calls after every block would defeat it. We have not observed this, and it
is the most obvious way the loop could be gamed.

\paragraph{Rules are searched on the rows they are scored on.} Every class
demotion, and the rules added from missed pushes, were chosen by reading this
corpus and reported on it. A session-wise split of the search transfers in
direction and barely in size: held-out precision rises in 16 of 20 splits, and
the mean gain is about two closings. The set of classes chosen is not stable
across halves.

\paragraph{The corpus is small where it counts.} 749 closings carry 82 pushes,
so every precision here has a Wilson interval two to four points wide and every
per-class number rests on tens of rows. No pooled configuration we measured
separates from the base rate at 95\%. Conclusions in this paper are stated at
the level the intervals support, which is rates and directions rather than
values.

\paragraph{Thresholds are tuned, not derived.} The demotion threshold and the
work threshold were chosen for the volume this log receives.

\section{Discussion}

If you run agents under standing permission, these three are the parts
we expect to travel. They are untested on a second user, and they were
measured on one corpus.

\paragraph{Put the decision where the capability is.} The clearest single result
is the background-work case: an exception that a 9B model could not hold in its
prompt became trivial once it was decided by ten lines of transcript inspection.
The temptation with a model in the loop is to route every judgement through
it. The cheaper lanes should take everything they can prove. On the closings
only the patterns catch, precision is 0.109, the base rate. On the closings
only the judge catches, it is 0.129. Among closings a person answered, those
unique judge rows reach 0.243 against 0.151. The model earns its place on the
rows a person saw, and none of it by being more discriminating than a regular
expression on the pooled table.

\paragraph{Confidence is a router before it is a decision.} Per-pattern
confidence is what makes the local model affordable, because it decides which
messages are worth a model at all. The judge's own confidence became usable
when the action space grew a third option: a STOP under the floor is put as
the question instead of as an accusation, and the turn is neither blocked nor
released. A weak signal that cannot pick between two actions can still pick a
third. Measured afterwards against the label, the low-confidence stops turn
out to be the most precise slice the judge produces, so the floor is keeping
its best rows and asking about them rather than discarding them. A confidence
read as a quality score would have inverted that.

\paragraph{Interventions leave labels.} Any system that interrupts an agent gets
a natural experiment for free: the agent's behaviour after the interruption. We
suspect this generalizes to most guardrails. The evaluation problem for
interventions is often not a missing label but an unread one.

\section{Limitations and future work}

This system is one developer's gate, scored on one machine. The next
three paragraphs are what that means for someone who wants to ship it.

The loop corrects unevenly. It finds patterns that block too much cheaply,
because every block leaves a graded trace and the report reads itself. The
other side needs the push label: a turn the gate allowed and the developer then
pushed back on is a miss, and reading those misses is still manual.
\texttt{STOP\_HOLDOUT} can suppress a fraction of fires at random to produce
clean miss-rate rows. The default is zero. A random miss on a sure
permission-ask is a failed re-wake, which is how the first Grok happy path
died one time in ten.

The scored tree cannot be graded on the frozen 749, because those
sessions never wrote \texttt{decisions.json}. The word leftover can,
and it fired on zero of them. A practitioner who wants this lane
to earn its keep has to write the tree while branching. The unit
fixtures grade the pick rule. They do not grade a year of live
trees.

Three extensions follow directly. Demotion is currently proposed and applied by
an agent reading the report; promotion in the other direction is not
implemented. The judge is a general instruction-tuned model, and the graded log
is exactly the dataset needed to fine-tune a specialist. And no part of this has
been run against a second user, which is the first thing that should happen. \texttt{SubagentStop} is a deliberate non-goal, not a missing port.

\section{Conclusion}

A gate that stops an agent from stopping early is easy to build and hard to
evaluate, because evaluation needs labels and labels need a reader. Both
labels this system needs were already lying in the transcript: the
developer's own reply, and the agent's own behaviour after the gate
intervenes. Neither costs an annotation.

On the stored file the gate fires on two closings in three and is right
about one in seven. On the closings a developer actually answered it is
right about one in five, while the failure runs at one in seven. The
association holds either way, $p = 0.0007$ pooled and $p = 0.0009$
stratified by week. 55 of 63 graded blocks bought work.

Against frozen 27B verdicts the live pattern file then takes person
recall from 0.829 to 1.000 on Claude and from 0.791 to 1.000 on Grok,
with zero extra false alarms. The two leftover extra-misses on Claude
are harness next-turns, not a person. A single fire is still a poor
bet: even at T4, four of five fires are false alarms.

A practitioner who ships this should leave it on, and should not trust a
single fire. Conversion is the number to run it on: after a wake, did
the agent work. Put the decision where the capability is. Do not scale
the model; scale around it. The next closing is still the one the gate
cannot name.

Neither the retired waiting exemption nor the prompt written as a test
came from a new idea. Both came from finally measuring something that
had been sitting in the transcript the whole time.
